% 第 11 章　转动的世界
\chapter{转动的世界}

\step{反常识开场}

\begin{mainline}
早餐桌上放着两个鸡蛋，一个生的，一个煮熟了。它们一样大、一样重、一样白，闭上眼睛摸不出任何区别。现在做一个三岁孩子都会的动作：把它们分别放在桌面上，用手指捻着转起来，像转陀螺那样。

结果立刻分出高下：熟鸡蛋转得又快又稳，像个精神抖擞的小陀螺，可以呼呼地转上十几秒；生鸡蛋却懒洋洋的，转不了两三圈就歪歪扭扭地停下来，好像有什么东西在蛋壳里面“拖后腿”。

这还不够怪。真正让人坐不住的是下一个动作：重新把熟鸡蛋转起来，然后用手指轻轻按住它——按大约一秒钟，让它完全停住——再松开。按理说这枚鸡蛋的转动已经被你“没收”了，可松开的一瞬间，它竟然自己又转了起来！是谁在里面重新拧了发条？鸡蛋里没有发动机，没有弹簧，只有蛋白和蛋黄。

先别急着翻答案。类似的怪事还有很多：花样滑冰运动员原地旋转，把张开的双臂猛地收拢，转速会凭空快上好几倍，没有人推她；自行车的轮子转起来之后就不容易倒，停下来却立刻东倒西歪；地球从形成至今自转了四十多亿年，从来没有人给它上发条，它却一天也没有忘记转。这些现象的背后，藏着一套和“推”同样基本、却常常被忽略的规律——转动的规律。本章的任务，就是把这套规律从鸡蛋、冰刀和自行车轮里挖出来。
\end{mainline}

\step{先猜一猜}

老规矩，先押注，再看牌。请认真写下你对下面三个问题的预测和理由，写错了不扣分——本章要修理的恰恰是直觉。

第一题。熟鸡蛋被你按住又松开后“复活”了。重新转起来的它，跟被按停之前相比，转速会（a）一样快；（b）明显变慢；（c）反而更快？

第二题。花样滑冰运动员收拢双臂，转速变成原来的三倍。她多出来的“快”是从哪来的？（a）冰面推了她一把；（b）她肌肉做了功，给自己加了速；（c）什么东西守恒了，重新分配的结果。

第三题。一扇关着的门，你用同样大小的力，分别尝试三种推法：推门把手（离门轴最远）、推门的正中间、贴着门轴旁边推。哪种推法最容易把门推开？如果你顺着门面、正对门轴的方向推门把手，又会怎样？

写下你的答案。等本章走完，我们会回来逐题对账。

\step{动手实验}

\begin{experiment}[鸡蛋转速表与“扭扭乐”木尺]
\textbf{材料}：一个生鸡蛋、一个熟鸡蛋（做好标记，别敲错了）、一张桌子；一把长尺或轻质木棍（半米到一米长）、两块一样大的橡皮泥或两团等重的湿纸巾、胶带。

\textbf{步骤一：辨蛋。}把两个鸡蛋分别横放在桌上，用差不多的手法捻转它们，各试三次。记录：哪个转得久？哪个转得稳？再对转得久的那个做“复活”实验：按住一秒，松开，观察它是否重新转起来，并比较复活前后的快慢。

\textbf{步骤二：光杆扭扭乐。}双手握住木尺正中间，让它水平，像拧毛巾那样快速地来回拧转（左转一点、右转一点，反复）。感受一下：让尺子改变转动状态，费劲吗？

\textbf{步骤三：戴上“哑铃”再拧。}把两块橡皮泥用胶带固定在尺子两端，再次握住中间来回拧转。同样的手、同样的拧法，费劲程度立刻不一样了——手臂明显地“带不动”。最后把橡皮泥挪到紧贴双手的位置（仍在尺上，总重量没变），再拧一次，又轻松了。

\textbf{记录}：尺子的总重量三次都一样，为什么“转起来、停下来、换方向”的难易天差地别？橡皮泥离手越远越费劲，这个“远”起了什么作用？连同鸡蛋的观察一起写下来，本章后面要逐条解释：你刚才亲手摸到的，是物理学里一个叫“转动惯量”的东西。
\end{experiment}

还有一个零成本实验，路过任何一扇门时都能做：用一根手指，推门把手处把门推开；再用一根手指，在门轴旁边用同样的力气推。第二次你会发现自己像个不会用门的婴儿。请记住这个感觉——“力一样大，效果不一样”，这就是本章全部故事的入口。

\step{旧模型遇到麻烦}

回顾我们手里的旧工具。第 6 章到第 10 章建立的模型非常成功：描述一个物体，报出它的位置和速度；预测它的未来，算清它受的合力，合力决定动量变化的快慢；碰撞前后动量守恒，过程前后能量守恒。这套语言管住了小车、台球和火箭。

可一遇到转动，麻烦就一个接一个地冒头。

麻烦一：力说不清转动。推门实验里，三次推门，力的大小完全一样、方向也一样，效果却天差地别。旧模型的公式 \(F=ma\) 里根本没有“力作用在哪儿”这个位置信息——它默认你推的是一个点。可门不是点，门是绕着门轴分布开来的一整块东西。要预言门会不会转、转多快，只知道“多大的力”不够，还必须知道“力落在离轴多远的地方”。旧模型缺了半条信息。

麻烦二：转动版本的“谁在续航”。第 6 章里，冰壶松手后一直滑，曾逼我们承认：匀速直线运动不需要谁来维持。现在同样的追问轮到转动：地球自转了四十多亿年，没有任何外力推它，它为什么不“累”？旧直觉会说“转着转着总会慢慢停吧”——可如果承认运动状态本身不需要解释、只有\emph{改变}才需要解释，那么转动状态也应当如此：没有什么去改变它，它就照原样转下去。旧模型里没有对应的概念来安放这句话。

麻烦三：鸡蛋的复活。按住熟鸡蛋时，它的转动明明停了一秒；松手后转动又回来了。第 10 章教我们：动量不会凭空产生，也不会凭空消失，只会转移。鸡蛋“复活”说明那一秒钟里，转动并没有真正消失，而是藏在了某个地方——藏在你手指按不住的蛋壳内部。可“转动藏起来了”是什么意思？旧模型只有“直线运动的账”（动量），没有“转动的账”。要听懂鸡蛋在说什么，我们得开一本新账。

三条麻烦指向同一个诊断：我们描述世界的语言里，缺了“转动”这一整套词汇。缺什么就造什么——这正是物理学几百年来的工作方式。

\step{发明新概念}

\begin{mainline}
\textbf{第一组概念：先学会描述转动——角位置、角速度、角加速度。}描述直线运动时，我们报位置、速度、加速度。描述转动，照方抓药：一个绕固定轴转动的物体，上面每一点到轴的距离不变，变的只是“转到哪儿了”。于是我们给转盘上某一条标记线报一个\emph{角位置}：它相对于起始方向转过了多少。角位置可以用“圈”来数，但更好用的单位是\emph{弧度}：标记点到转轴的距离是 \(r\)，它沿圆周走过的弧长除以 \(r\)，就是转过的弧度数。这样规定的好处立刻显现：弧长等于“半径乘以弧度”，距离轴两倍的点，同样的转角走两倍的路——一个数就管住了半径上所有的点。转一整圈，弧长是整个圆周，对应的角是 \(2\pi\) 弧度，约 \(6.28\)。

角位置随时间变化的快慢，叫\emph{角速度}，意思是“每秒转过多少弧度”（日常生活中“每分钟多少转”是它的亲戚）；角速度本身变化的快慢，叫\emph{角加速度}。这三个量是位置、速度、加速度在转动世界的直系亲属，家族的姓都一样：一个管“此刻在哪”，一个管“正在变多快”，一个管“变化本身变多快”。

墙上的挂钟是个现成的角速度陈列柜：时针、分针、秒针装在同一根轴上，角速度却各不相同——秒针每分钟一圈，分针每小时一圈，时针半天才一圈，三者之比是 \(60:1:\tfrac1{12}\)。更有意思的是同一根秒针：轴心附近的点和针尖的\emph{角速度}完全一样（一分钟都转一整圈），\emph{线速度}却差了几十倍，针尖每分钟要跑过将近十厘米的弧长，轴心几乎原地不动。记住这个区别：转动世界里“大家共享的是角速度，各跑各的是线速度”，后面很多错觉都是把这两者混淆造成的。

\textbf{第二个概念：力矩——力的“转动效果”不是由力单独决定的。}推门和拧螺丝的经验可以浓缩成一句话：一个力能产生多大的转动效果，取决于两个因素的乘积——力有多大，以及转轴到这个力的作用线的垂直距离（叫\emph{力臂}）有多长。我们把这个乘积叫\emph{力矩}。力矩大，转动状态就改变得快；力矩为零，转动状态就不变。它像什么？它像用杠杆撬东西时的“劲”：力是原料，力臂是放大倍数。它不像什么？它\emph{不是}一种新力——门轴附近并没有藏着什么神秘的东西，力矩只是“力加位置”合伙产生的效果，就像第 9 章里“功”是“力加距离”合伙的效果一样。这解释了为什么门把手总装在离轴最远的地方、扳手要做成长柄、拧不动的螺母套上钢管就能拧动：工程上从不浪费力臂。

环顾四周，力矩的设计无处不在。螺丝刀的手柄做得粗，是为了让你的握力落在离轴更远的地方——同样一把螺丝刀，换粗柄的立刻好拧；汽车方向盘比游戏方向盘大得多，大巴的方向盘又比小轿车的大，都是同一个道理；开罐头时把勺柄插进拉环底下撬，用的是罐头边缘做支点；连你拧干毛巾时，都会下意识地把手张开到毛巾两端，而不是挤在中间。古人早就把这条规律用到了极致：秤杆上一个小小的秤砣，靠挪动位置就能称出几十斤的货物——秤星标的就是力臂，不是重量。
\end{mainline}

\begin{center}
\begin{tikzpicture}[scale=1.0]
  % 俯视图：门（从门轴向右展开）
  \fill[pattern=north east lines] (-0.15,-0.9) rectangle (0.15,0.9);
  \node[left] at (-0.2,0) {门轴};
  \draw[very thick,gray!60!black] (0,0) -- (5.6,0);
  \node[above] at (4.6,0.02) {门（俯视）};
  % 力1：推门把手，垂直于门
  \draw[-{Stealth[length=3mm]},very thick,red!70!black] (5.2,0) -- (5.2,1.6) node[above] {力 \(\to\) 力臂最长：效果最好};
  \draw[fill] (5.2,0) circle (0.05);
  % 力2：推中间
  \draw[-{Stealth[length=3mm]},very thick,orange!80!black] (2.8,0) -- (2.8,1.6) node[above,yshift=16pt] {同样的力，力臂减半：效果减半};
  \draw[fill] (2.8,0) circle (0.05);
  % 力3：贴着轴推
  \draw[-{Stealth[length=3mm]},very thick,blue!70!black] (0.9,0) -- (0.9,-1.3) node[below] {力臂几乎为零：推不动};
  \draw[fill] (0.9,0) circle (0.05);
\end{tikzpicture}
\end{center}

\begin{mainline}
\textbf{第三个概念：转动惯量——抗拒转动改变的，不只是质量，还有质量的“站位”。}“扭扭乐”实验里，木尺加橡皮泥的总重量没变，改变的是重量离转轴的远近：离轴越远，来回拧转越费劲。这说明“对转动改变的抗拒”不只取决于有多少质量，还取决于质量分布在哪儿。我们给这个抗拒程度起名\emph{转动惯量}。它是质量（第 6 章里抗拒速度改变的程度）在转动世界的对应物，但多了一条脾气：\emph{离轴越远的质量，“发言权”越大}，而且是按距离的平方放大——同一块橡皮泥挪到两倍远，抗拒变成四倍。所以花样滑冰运动员张开双臂时转动惯量大、转速慢，收拢双臂时转动惯量骤减；所以超市买的实心小月饼盒和同样重量的空心大铁环，转起来完全是两种脾气。

玩具设计师把这条脾气用得很彻底。一只好的悠悠球，尽量把重量做在轮盘的\emph{外缘}而不是轴心附近，转动惯量做大后，它就能以不高的转速携带可观的角动量，在绳子尽头长时间“睡眠”，经得起各种花式动作；街边抽的陀螺则是反过来利用这一点——鞭子每抽一下，都是一次切向的力施加在远离轴的边缘，力矩大，陀螺越转越快，而圆滚滚的重心配置让它的质量分布匀称，转起来稳如小树桩。还有杂技演员头顶转的长杆：杆顶绑上一只篮子或盘子，故意把质量放高放远，转动惯量一大，杆倒下来的角加速度就小，演员才有充足的时间挪步去救——转得慢，在这里是救命的本事。

\textbf{第四个概念：角动量——转动世界的“运动量”。}第 10 章里，我们把“运动的量”定义为动量，并发现它在碰撞中守恒。转动世界有对应的量：\emph{角动量}，等于转动惯量乘以角速度。转动惯量回答“这家伙多难被拧动”，角速度回答“它现在转多快”，乘起来就是“它此刻携带着多少转动”。而转动世界最深的规律是：\emph{一个系统所受的外力矩加起来为零时，它的总角动量保持不变}。这就是角动量守恒，是动量守恒的转动版，可靠程度一模一样：直到今天，没有任何实验发现过例外。

\textbf{一张对照表，把两本账对齐。}至此，转动世界的每个概念都能在平动世界找到亲戚。下面这张表值得抄在笔记本上——本章剩下的内容，包括以后遇到的所有转动问题，几乎都是在这张表的左右两列之间做翻译。
\end{mainline}

\begin{center}
\begin{tabular}{lll}
\toprule
\textbf{角色} & \textbf{平动世界} & \textbf{转动世界} \\
\midrule
状态：此刻在哪儿 & 位置 & 角位置（弧度） \\
状态：正在变多快 & 速度 & 角速度 \\
状态：变化在变多快 & 加速度 & 角加速度 \\
抗拒改变的程度 & 质量 & 转动惯量 \\
改变运动状态的能力 & 力 & 力矩 \\
“运动携带的量” & 动量 & 角动量 \\
运动状态改变的规律 & 合力决定动量变化 & 外力矩决定角动量变化 \\
守恒律 & 无外力：动量守恒 & 无外力矩：角动量守恒 \\
\bottomrule
\end{tabular}
\end{center}

\step{一句话模型}

\begin{oneline}
改变转动状态的不是力本身，而是力矩（力乘力臂）；转动惯量衡量“多难被拧动”，质量离轴越远越大；只要外力矩加起来为零，系统的角动量——转动惯量乘角速度——就保持不变，它只会在系统内部重新分配。
\end{oneline}

背下这一句话，你就同时拿到了解释鸡蛋、滑冰、陀螺和中子星的钥匙。

\step{数学翻译器}

\begin{mainline}
先翻译描述部分。角位置 \(\theta\)（弧度）、角速度 \(\omega\)、角加速度 \(\alpha\) 的定义，和直线运动一模一样，只是把“路程”换成“转角”：
\[
  \omega=\frac{\Delta\theta}{\Delta t},\qquad
  \alpha=\frac{\Delta\omega}{\Delta t}.
\]
读法：\(\omega\) 回答“每秒转多少弧度”，\(\alpha\) 回答“角速度每秒变多少”。例行检查：若 \(\omega=0\)，物体不转；若 \(\omega\) 很大而 \(\alpha=0\)，物体飞快但\emph{匀速}地转——地球就是这样；\(\alpha\) 为负不一定表示“越转越慢”，要看它跟 \(\omega\) 同号还是异号，这和直线运动里“加速度为负不等于减速”是同一个坑。

接着翻译力矩。力臂为 \(d\)（转轴到力的作用线的垂直距离）、力为 \(F\) 时：
\[
  \tau=Fd .
\]
三道例行检查。第一道，\(d=0\)：力穿过转轴，力矩为零——贴着门轴推门、顺着门面朝门轴推，都推不开门，正是先猜一猜第三题的后半问。第二道，力臂加倍，\(\tau\) 加倍：这解释了为什么拧生锈的螺母要换长扳手，也解释了天平为什么“小砝码挪远格能顶大砝码”。第三道，方向：力矩有正负（顺时针与逆时针），两个力矩互相抵消时，物体的转动状态不变——推门的人若在门的另一侧用相反的力矩推，门就僵持不动。

代入真实数字感受一下。修自行车时拧一颗脚踏板螺母，规定拧紧力矩约 \(30\ \mathrm{N\cdot m}\)。用一把 25 厘米长的扳手，你需要施加 \(30/0.25=120\) 牛的力，相当于提起一桶 12 千克的水；手劲不够怎么办？踩上去——把 60 千克体重的一部分压到扳手末端，力臂不变，力翻几倍，力矩轻松达标。反过来，汽车维修手册要求某些精密螺栓只能拧到 \(10\ \mathrm{N\cdot m}\)，这时老司机会\emph{握着扳手的中段}拧，主动缩短力臂，防止手滑拧断螺栓。同一个公式，正着用是放大力量，反着用是限制力量。

再翻译转动惯量。想象物体由许多小块组成，第 \(i\) 小块质量 \(m_i\)、离轴距离 \(r_i\)，则
\[
  I=\sum_i m_i r_i^2 .
\]
这是全章唯一一个“天生带平方”的式子，值得盯它一眼：距离的平方意味着把质量挪远是\emph{最昂贵}的增惯手段。检查：所有质量都在轴上（\(r_i=0\)），\(I=0\)，这根“理想细杆”绕轴转动毫不费力；同一块质量挪到两倍远，贡献变四倍——“扭扭乐”实验的手臂感受被精确复现。两个常用结果（推导见数学桥层）：质量 \(M\)、半径 \(R\) 的\emph{薄圆环}绕中心轴，所有质量都在距离 \(R\) 处，\(I=MR^2\)；同质量同半径的\emph{实心圆盘}，质量平均来说离轴近得多，\(I=\tfrac12 MR^2\)。一样重、一样大的环和盘，环的转动惯量是盘的两倍——这个差别会在本章挑战题里滚下斜坡。

\begin{center}
\begin{tikzpicture}[scale=0.85]
  % 薄圆环：质量全在边缘
  \draw[very thick,red!70!black,fill=red!8] (0,0) circle (1.5);
  \draw[very thick,red!70!black,fill=white] (0,0) circle (1.25);
  \draw[-{Stealth[length=2.5mm]},thin] (0,0) -- (1.06,1.06) node[midway,above left=-2pt] {\(R\)};
  \node at (0,-1.95) {薄圆环：\(I=MR^{2}\)};
  % 实心圆盘：质量铺满
  \begin{scope}[xshift=5.2cm]
    \draw[very thick,blue!70!black,fill=blue!10] (0,0) circle (1.5);
    \foreach \r in {0.45,0.9} \draw[thin,blue!40!black] (0,0) circle (\r);
    \draw[-{Stealth[length=2.5mm]},thin] (0,0) -- (1.06,1.06) node[midway,above left=-2pt] {\(R\)};
    \node at (0,-1.95) {实心圆盘：\(I=\tfrac12 MR^{2}\)};
  \end{scope}
\end{tikzpicture}
\end{center}

最后翻译主角。角动量
\[
  L=I\omega ,
\]
守恒律说：外力矩之和为零时，\(L\) 不变。写成前后对照就是 \(I_1\omega_1=I_2\omega_2\)。检查极限：\(I_2\) 减半，\(\omega_2\) 翻倍——转动惯量和角速度像跷跷板的两头，一个下去，另一个必须上来。
\end{mainline}

\textbf{一笔真实数据的估算：花样滑冰。}一位运动员旋转时，伸开双臂和一条腿，转动惯量约 \(3.2\ \mathrm{kg\cdot m^2}\)（体重约 60 千克，大部分质量分布在半径二三十厘米处），以每秒 1 圈起转，即 \(\omega_1=2\pi\ \mathrm{rad/s}\)，角动量约 \(3.2\times6.28\approx20\ \mathrm{kg\cdot m^2/s}\)。收拢双臂贴紧身体后，质量收缩到半径十厘米以内，转动惯量降到约 \(1.0\ \mathrm{kg\cdot m^2}\)。冰面摩擦力矩很小，几秒内角动量近似守恒，于是
\[
  \omega_2=\frac{I_1}{I_2}\omega_1\approx\frac{3.2}{1.0}\times 1\ \text{圈/秒}\approx 3.2\ \text{圈/秒}.
\]
这正是比赛录像里“收臂瞬间转速暴涨”的数量级——三周半、四周转体跳赖以成立的物理。顺带一笔账：收臂后转动动能 \(E=\tfrac12 I\omega^2\) 从约 \(63\) 焦升到约 \(200\) 焦，多出的能量不是守恒律白送的，而是运动员肌肉把手臂拉近身体时对抗“甩出去的劲儿”所做的功——角动量守恒和能量守恒各管各的账，第 9、10 章学到的分工在这里依然有效。

\begin{mathbridge}
为什么用弧度而不是“度”？因为只有用弧度，弧长公式才干净：离轴 \(r\) 处的点，线速度 \(v=\omega r\)，线加速度的切向部分 \(a_{\text{切}}=\alpha r\)。注意这个公式的推论：同一个转盘上，\(\omega\) 处处相同，\(v\) 却随 \(r\) 正比增长——转盘边缘的点比中心附近的点“跑”得快得多，尽管它们角速度完全一样。这正是“转动”和“平动”最不一样的地方：\(\omega\) 是整个刚体共享的一个数，\(v\) 却各点不同。

力矩的完整写法是 \(\tau=rF\sin\varphi\)，其中 \(\varphi\) 是位矢与力的夹角。它自动包含两个极端：力穿过轴心（\(\varphi=0\) 或 \(\pi\)）时力矩为零；力垂直于力臂时力矩最大。写成矢量，\(\bm{\tau}=\bm r\times\bm F\)，方向由右手定则给出，沿转轴；同样地，\(\bm L=I\bm\omega\) 也是沿轴的矢量。转动的牛顿第二定律写作
\[
  \bm\tau_{\text{外}}=\frac{d\bm L}{dt},
\]
平动与转动的对照在这行式子里达到顶峰：只要把 \(\bm F\) 换成 \(\bm\tau\)、\(\bm p\) 换成 \(\bm L\)，形式一字不差。圆盘转动惯量 \(I=\tfrac12MR^2\) 的来历也只是一次积分：把圆盘切成无数细环，半径 \(r\) 处的环贡献 \(r^2\,dm\)，从 \(0\) 积到 \(R\) 即得。
\end{mathbridge}

\begin{challenge}
陀螺为什么不倒，反而会“转着圈点头”？简化版答案：重力对陀螺支点施加的力矩 \(\bm\tau\) 不为零，按 \(\bm\tau=d\bm L/dt\)，它不断改变角动量\emph{矢量}的方向而不是大小——\(\bm L\) 的箭头被迫绕竖直方向匀速扫出一个圆锥，这就是进动。严格处理需要把刚体的角速度分解到三个惯量主轴上（欧拉方程），那时会发现自由转动的刚体绕最大或最小惯量轴稳定，绕中间轴则会发生著名的“翻转不稳定”——在国际空间站的舱内抛转一把 T 形扳手，就能看到它周期性地翻身。这部分需要矢量微积分，本书不做要求，但请记住结论的形状：\emph{力矩改变的是角动量这个箭头的指向}，这是自行车轮、陀螺仪和回转罗盘全部本领的根源。
\end{challenge}

\step{模型的边界}

每个模型都要交代自己的势力范围，转动模型有四条边界最容易踩线。

\textbf{边界一：刚体假设。}本章的公式默认研究对象是\emph{刚体}——形状固定、各部分一起转的物体。熟鸡蛋近似满足，生鸡蛋完全不满足：蛋壳一转，里面的蛋清蛋黄靠黏糊糊的摩擦才慢慢跟上。一旦物体内部可以相对滑动，就不能再用一个 \(\omega\) 描述它，而要分块记账。猫从空中翻身、跳水运动员空中收展身体，都是“非刚体”的高级玩法，原理仍是角动量守恒，但计算远为复杂。

\textbf{边界二：守恒的条件是“外力矩为零”，不是“外力为零”。}两个最容易混。一个物体可以受着不小的合力却角动量守恒（比如被轴光滑固定住的飞轮，轴的支持力很大，但力臂为零，力矩为零）；也可以合力为零却越转越快——门两侧各站一人，用大小相等、方向相反的力同时推同一侧的门把手两端？更标准的例子是方向盘：双手一上一下反向用力，合力为零，力矩不为零，方向盘照样转起来。判断转动，永远先问力矩。

\textbf{边界三：定轴与自由轴。}我们默认转轴固定。轴一旦自由，事情会丰富得多：陀螺进动、地球自转轴两万六千年的缓慢画圈（岁差）、抛出的扳手翻身。现象惊人，账本的原理没变，变的只是“哪个方向算守恒”的几何。

\textbf{边界四：别给旧误解续命。}两条要专门点名。其一，转弯时“把人往外甩的离心力”不是一种新的力：让汽车转弯的向心力由轮胎与地面的摩擦提供，让转动雨伞上水珠沿切线飞出去的也不是“被甩”，而是水珠失去了继续做圆周运动所需的牵拉——向心力和力矩一样，都是对已有力（重力、弹力、摩擦）在某个问题里的\emph{角色描述}，不是与它们并列的新成员。其二，一个流传甚广的过度解释值得当着反例拆掉：“自行车不倒全靠车轮的陀螺效应”。陀螺效应确实存在，但 2011 年有研究者造出了一种特殊自行车——前后各装一对反向旋转的轮子，把整车的转动角动量精确抵消为零——它骑行时照样能自我稳定。真正让自行车稳的，是前轮转轴的几何设计（前叉后倾，使前轮自动向倾斜方向偏转）加上骑车人不停的微调。这个反例是本章最该记住的一课：\emph{一个正确的物理机制，不等于它是唯一在起作用的机制}。手机里的陀螺仪芯片、卫星上的姿态控制飞轮倒是货真价实地吃着角动量这碗饭：卫星不带燃料就能转身，靠的是舱内几只电动机驱动的飞轮——飞轮加速正转，星体就缓缓反转，角动量在二者之间转移，总量纹丝不动，这叫“动量轮”，哈勃望远镜就是用它对准深空目标的。

\step{回头解释开场现象}

现在回到早餐桌上那两只鸡蛋，用新模型逐条对账。

生鸡蛋为什么转不动？你的手捻动的是\emph{蛋壳}，蛋壳获得角速度，但里面的蛋清蛋黄是液体，只靠黏滞摩擦被蛋壳拖着走。结果是：壳转得比内容物快，内外角速度不一致——这不是刚体，边界一被触发。内外之间的黏性力矩不断把角动量从壳向液体转移，同时桌面摩擦持续从整个系统抽走角动量，于是整体很快慢下来。你付出的那一份角动量，大半耗散成了蛋液内部微乎其微的热。

熟鸡蛋为什么转得久？煮熟后蛋白蛋黄凝固成一个整体，近似刚体：手一捻，内外一起获得角速度，没有内耗，只剩桌面和空气的微小阻力矩，角动量被它抱得紧紧的，自然转得又稳又久。

最精彩的是“复活”。按住熟鸡蛋的那一秒，你的手指通过摩擦让\emph{蛋壳}停下；但凝固的蛋黄和蛋白并非理想刚体，尤其是蛋黄，它与蛋白之间仍有微小的可滑动余地——按住的瞬间，内部还残留着一部分没有停下来的转动，角动量并没有全交给你的手指，一部分藏进了蛋内。松手后，内部的黏性摩擦反向起作用：还在转的内容物通过内摩擦力矩重新拖动蛋壳，整个蛋又转起来。账本完美闭合：复活后的转速\emph{一定比原来慢}，因为按住期间手指和桌面已经抽走了一部分角动量（先猜一猜第一题答案：b）。如果你观察得够细，会发现按得越久，复活越弱——按到内部也彻底停住，就再也复活不了了。

花样滑冰的账本章已算过：收臂把质量从远处拉回轴边，转动惯量降到三分之一，角动量守恒逼角速度涨到三倍。自行车轮的谜底则是“陀螺效应 \(+\) 前叉几何 \(+\) 人的微调”的合奏——三个声部缺一不可，把任何单一机制吹成全部原因，都是对账本的误读。

最后轮到开场里的第三个悬念：地球为什么转了四十多亿年还不停？现在答案几乎是白送的——地球近似一个漂浮在真空里的巨大刚体，太阳和月球的引力几乎不对它施加净力矩（引力作用线穿过地心附近），于是它的自转角动量无处流失，只能保持。它不是“不知疲倦”，而是根本不需要力气：转动状态和匀速直线运动一样，是不需要解释、也不需要维持的默认状态。当然，账要算细：月球掀起的潮汐与海底摩擦，像一只极轻柔的手按在地球这颗“熟鸡蛋”上，每世纪从地球自转里抽走一点点角动量，一天因此每世纪变长约 \(1.7\) 毫秒——四十多亿年里，一天已经从当初的几个小时拉长到二十四小时。这笔微小的外流去了哪里，正是本章挑战 2 要追的账。

\step{脑洞实验与挑战题}

\begin{thought}[悬空的宇航员能不能转身？]
一名宇航员飘在太空舱正中央，四肢碰不到任何舱壁。她浮力中性，纹丝不动。她能\emph{只靠自己}改变朝向、从面对舱门变成背对舱门吗？

直觉说不能：没有支点，没有外力，总角动量是零而且必须保持为零，怎么转身？但她可以。做法（猫在下落时用的就是这套）：先把双臂向左侧平伸，上半身向右拧转——因为总角动量必须为零，下半身和腿会反向（向左）转一个角度；然后收拢双臂、把腿弓起换个姿势，再反向拧回。由于收拢和伸展时上下半身的转动惯量不同，两次拧转产生的角位移不能互相抵消，一套动作下来，总角动量始终是零，而她的朝向已经改变。这不是魔术，是“非刚体的角动量守恒”：守恒限制的是\emph{总量}，不是\emph{姿态}。太空站的宇航员确实用类似动作调整身体方向——零外力矩的世界里，转动的规则不是“不许动”，而是“每一分转动都必须有反向的转动作陪”。
\end{thought}

把尺度推到宇宙的另一头，角动量守恒会写出更狂野的故事。一颗质量与太阳相当的恒星，晚年核心坍缩：半径从约 \(7\times10^5\) 千米（太阳现在的半径）缩到约 \(10\) 千米，几乎全部质量保留。转动惯量正比于半径的平方，于是 \(I\) 缩到原来的约 \((10/7\times10^5)^2\approx2\times10^{-10}\)；若坍缩过程角动量大致守恒，自转角速度就要放大约五十亿倍。太阳自转一圈约 25 天，坍缩后的核心——一颗\emph{中子星}——理论上每秒可转上千圈。真实观测到的毫秒脉冲星，最快的每秒自转超过 700 圈，与这个估算同一数量级：赤道表面的线速度逼近光速的几分之一，而它只是把“花滑运动员收臂”演到了极致。从早餐的鸡蛋到恒星尸体，管账的是同一条定律。

\begin{puzzles}
\item 同一斜面顶端，同时放手一个实心圆盘和一个薄圆环（质量、半径相同，都只滚不滑）。谁先滚到底？提示：重力势能不变，全部转化成动能，但动能要分成“平动”和“转动”两份；转动惯量大的，转动那份“吃”得多，留给平动的就少。回忆 \(I_{\text{环}}=MR^2\) 与 \(I_{\text{盘}}=\tfrac12MR^2\)，答案与质量、半径都无关。
\item 天文测量发现：地球上的一天每世纪约变长 \(1.7\) 毫秒，同时月球以每年约 \(3.8\) 厘米的速度远离地球。这两件事为何总是一起出现？提示：潮汐摩擦像一只手按在地球这个“熟鸡蛋”上；地月系统的外力矩近似为零，地球自转丢掉的角动量去了哪里？月亮变远，它的轨道角动量是变大了还是变小了？
\item 冰上旋转的运动员收臂加速后，把手套沿\emph{切线方向}抛出去。抛出的瞬间，她的转速会立刻大变吗？提示：手套沿切线飞出，它带走的角动量与它的运动方式有关；先想清楚“角动量守恒”管的是谁——是她一个人，还是“她加手套”？再想想，如果她不是抛出而是\emph{轻轻放手}，区别在哪。
\end{puzzles}

本章的两条延伸线，先埋在这里。其一，角动量守恒这条铁律本身将在第 58 章被“追究来历”：它不是碰巧成立的经验，而是“空间的各个方向等价”这一条对称性的直接产物——这是二十世纪物理最深刻的发现之一，守恒律统统来自对称性。其二，到了微观世界，角动量会露出它的量子面孔：电子、质子携带着一种名为“自旋”的内禀角动量，它不能理解为小球真的在自转，数值还是一份一份量子化的——那是第 47 章的故事。但无论你走到多深，账本的格式都不会变：先问系统此刻的转动状态，再问什么力矩在改变它，最后问我们凭什么测到它——转动世界和平动世界，用的是同一套追问。
